Geometric classes are represents objects in three-dimensional space. These include points, directions, and 3D volumes. Points are implemented as \lstinline|Eigen::Vector3d| objects (the \lstinline|3d| here means the
vector contains 3 \lstinline|double|s). Directions are represented by the \lstinline|UnitVector3d| class, which encapsulates a normalized \lstinline|Eigen::Vector3d| member. In this section, \lstinline|Point| is an alias for
\lstinline|Eigen::Vector3d| and \lstinline|Direction| is an alias for \lstinline|UnitVector3d|.

3D volumes are derived from the abstract \lstinline|Shape| class. They must have the following functions overridden in order to be instantiated:
\begin{enumerate}
    \item \lstinline|xMin()|, \lstinline|xMax()|, and so on: locate the extreme points in each of the three Cartesian coordinates.
    \item \lstinline|surfaceArea()|
    \item \lstinline|volume()|
    \item \lstinline|surfaceContains(const Point&)|: checks if the Point is strictly on the surface.
    \item \lstinline|encloses(const Point&)|: checks if the point is strictly inside \lstinline|*this| (the \lstinline|Shape| in question). This means the point is not on the surface or outside.
    \item \lstinline|overlaps(const Shape&)|: checks if the two shapes have a positive volume of intersection. That means if the intersection is a point, a line, or a plane, it returns \lstinline|false|.
    \item \lstinline|distanceToSurface(const Point&, const Direction&)|: calculates the smallest, positive distance that a point has to travel to reach the surface.
    \begin{itemize}
        \item If the point is inside the shape, the distance equation will have positive and negative roots, and only the positive root is returned.
        \item If the point is on the surface, returns 0 if it leaves the surface, and the distance to the other end if it enters the surface.
        \item If the point is outside the shape, returns NAN if it never enters the shape, else the smallest distance to the surface.
    \end{itemize}
    \item \lstinline|normal(const Point&)|: checks if the point is on the surface, and returns the outward normal vector if it is.
    \item \lstinline|print(std::ostream& os)|: outputs the shape into \lstinline|std::cout|.
\end{enumerate}

\lstinline|Shape| also has an overloaded \lstinline|encloses(const Shape&)| that checks if the other shape is completely contained within \lstinline|*this|. This function evaluates to \lstinline|true| even if
the two \lstinline|Shape|s intersect on the surface. Because of the tree design, this is made a function within \lstinline|Shape|, but will always throw an exception unless \lstinline|*this| is of type
\lstinline|Box|, and practically enclosure check is always done on \lstinline|Box| objects.